encloses a period-28 ring around a chaotic core. Period-21 windows sit in the chaotic region as thin crescents with
legs, and curves of very negative exponent, where the orbits are superstable (the leading multiplier crosses zero),
run through the windows as the ``endoskeleton'' of shrimps does \cite{facanha2013endoskeletons}. In this zoom
18\,\% of cells are chaotic.

\paragraph{No one-dimensional shrimp centre.} In a shrimp of a one-dimensional map, two superstable curves of the
window's orbit cross at a doubly superstable centre \cite{gallas1993structure,facanha2013endoskeletons}. We continued
one period-7 orbit over a $31\times31$ grid of the body's upper part and followed the sign of its leading multiplier
(Figure~\ref{fig:superstable}). Its superstable curves form two arcs that approach within $1.3\times10^{-3}$ in
$\log_{10}T$ near $(\alpha,\log_{10}T)=(0.0700,0.206)$ and do not cross at this resolution. A five times finer grid of that approach
(1155 cells) shows why. There the leading multiplier is tiny (median modulus $9\times10^{-4}$) and as small as the next
one, and in 221 cells the two form a complex pair, in two wedges that flank the gap. On this sampled approach the two
computed leading multipliers exchange dominance and form a complex pair, so the sign of a leading real multiplier no
longer defines a scalar superstable curve there. Where a one-dimensional map would place its doubly superstable centre,
the orbit's local linearization has two comparable directions. Next to the body, the 7-step return maps of chaotic
cells are single-humped in every phase, with doubled strands in two phases, the trace of a thin second dimension
(Figure~\ref{fig:returnmapsbody}); cobwebs on these maps show the doubling cascade and the chaos as on a
one-humped interval map (Figure~\ref{fig:cobwebs}). On the tracked branch and in the sampled region, the body is
therefore not the shrimp of a one-dimensional map. Whether it is the shrimp of an effectively
two-dimensional map, as in the H\'enon family, is open.

\paragraph{What the shrimp-like shapes are.} The shapes are made of the bifurcation curves of single windows
(Figure~\ref{fig:shrimpschematic}). A periodic window is the region where one periodic orbit attracts. The windows
we continued are bounded by fold curves, where the orbit is born or dies together with a saddle (multiplier $+1$), and
flip curves, where it doubles (multiplier $-1$); in higher dimensions other boundaries, such as a complex pair crossing
the unit circle, are also possible. In a one-dimensional map the multiplier runs from $+1$ through $0$ to $-1$ across the
window, so every window holds a superstable curve: the dark ``endoskeleton'' lines of our maps. In a higher-dimensional
map the leading multiplier can instead change sign through a complex pair, as it does in parts of our windows. Beyond
some of the flip curves we followed lie the doubling sequences of Section~\ref{sec:cascade}; a flip need not start a
complete cascade.

Steering strength and temperature move the orbit in different directions, so the fold curves bend. Where a fold curve
turns back, a cusp can form (a visual turn alone is not our criterion: the located cusps satisfy the degeneracy and
unfolding checks of Section~\ref{sec:curves}), and between the two branches leaving the cusp the window continues as a
thin leg.
Bent fold curves with cusps, flip curves inside them and a superstable line through the body produce the bodies, legs,
hooks and crescents of the maps. All nine reversals of the fold curves we continued are located cusps
(Section~\ref{sec:curves}).

A shrimp in the sense of Gallas has one more element: a centre, where two superstable curves of the same orbit cross.
There the orbit passes two critical points at once, which needs two turning points in a one-dimensional reduction.
The return maps of the hostile body have one turning point per phase (Figure~\ref{fig:returnmapsbody}). Where the hostile body's superstable curves approach each other, the orbit becomes
two-dimensional instead of doubly superstable (Figure~\ref{fig:superstable}). The period-6 hook's orbit never comes
close to superstable: its leading multiplier stays above $0.028$ in modulus.

The X-shaped crossing of a period-5 and a period-6 strip (Figure~\ref{fig:zooms}, bottom) is not one window's anatomy
but two windows overlapping, and at the crossing both of their orbits are unstable (Section~\ref{sec:curves}). The
shapes are therefore the two-parameter unfolding of single-humped windows. They are shrimp-like in body, legs and
internal cascade, but have no shrimp centre.
\input{figures/tikz-shrimp}
\begin{figure}[t]
\centering
\includegraphics[width=\linewidth]{fig-superstable-lines.png}
\caption{Superstable structure of the hostile period-7 body (Qwen2.5-0.5B-Instruct, hostility direction at layer 12,
float64). (a) Largest Lyapunov exponent of the period-7 cells, with fold (black) and flip (red, dashed) curves and the
two located cusps (gold); dark curves are superstable lines. (b) One period-7 orbit continued over the red box of (a):
sign of its leading multiplier (red positive, blue negative), dark where it is small, and its zero curves (black).
The two arcs approach near $(0.0700,0.206)$ without crossing. (c) Five times finer zoom of the cyan box of (b):
logarithm of the leading multiplier's modulus. Crosses mark cells where the two leading multipliers form a complex
pair; white segments mark sign changes between cells with a real leading multiplier. Where a one-dimensional map would
have its centre, the orbit is two-dimensional.}
\label{fig:superstable}
\end{figure}
\begin{figure}[t]
\centering
\includegraphics[width=\linewidth]{fig-return-maps-body.png}
\caption{Seven-step return maps $s_{n+7}$ against $s_n$, for each phase of the 7-periodic text, of chaotic cells next to
the hostile period-7 body (top three rows: $\alpha=0.07136$--$0.07142$, $\log_{10}T=0.2057$--$0.2062$) and on the
hostile line of Figure~\ref{fig:hubcascade} (bottom row: $\alpha=0.0655$, $\log_{10}T=0.1848$); $s$ is the projection on
the leading principal component of that phase's states, which carries 89--100\,\% of their variance. Every phase
shows a single hump. In phases 4 and 5 the curve splits into two thin strands.}
\label{fig:returnmapsbody}
\end{figure}
